\documentclass[10pt,a4paper]{amsart}
\usepackage[margin=19mm]{geometry}
\usepackage{amsmath,amssymb,mathtools}
\usepackage[scr=rsfs,cal=euler]{mathalfa}
\usepackage{xfrac}
\usepackage[unicode,hidelinks,bookmarks=false]{hyperref}
\newcommand{\ie}{i.e.\ }
\newcommand{\cf}{cf.\ }
\newcommand{\resp}{respectively\ }
\newcommand{\Subsection}{\subsection}
\providecommand{\nobreakdash}{\nobreak\hbox{-}}
\setlength{\emergencystretch}{2em}
\multlinegap0pt
\usepackage{bm}
\usepackage{bbm}
\usepackage{dsfont}
\usepackage{xspace}
\usepackage{cancel}
\usepackage{mathtools}
\usepackage{cleveref}
\usepackage{amsthm}
\makeatletter
\@ifundefined{lemma}{\newtheorem{lemma}{Lemma}}{}
\@ifundefined{theorem}{\newtheorem{theorem}{Theorem}}{}
\makeatother
\newcommand{\R}{\mathbb{R}}
\newcommand{\Ttran}{\mathsf{T}}
\newcommand{\T}{\mathsf{T}}
\newcommand{\defi}{:=}
\newcommand{\range}{\mathsf{range}}
\newcommand{\tol}{\mathsf{tol}}
\providecommand{\Bignorm}[1]{\Bigl\lVert#1\Bigr\rVert}
\providecommand{\abs}[1]{\lvert#1\rvert}
\providecommand{\bignorm}[1]{\bigl\lVert#1\bigr\rVert}
\providecommand{\diagm}{\mathrm{diag}}
\providecommand{\norm}[1]{\lVert#1\rVert}
\providecommand{\spa}[1]{\mathsf{span}\{#1\}}
\title[Lanczos with compression for symmetric eigenvalue problems]{Lanczos with compression for symmetric eigenvalue problems}
\author{Angelo A. Casulli, Daniel Kressner, Nian Shao}
\date{Source: arXiv:2602.20948v2 (CC BY 4.0)}
\begin{document}
\maketitle

\noindent\emph{Source and licence.} Lanczos with compression for symmetric eigenvalue problems. Version arXiv:2602.20948v2 (CC BY 4.0), distributed under the Creative Commons Attribution 4.0 International licence, as recorded in the arXiv licence metadata for this exact version and shown on the abstract page https://arxiv.org/abs/2602.20948v2. Full rights notice: licenses/arxiv-2602.20948-CC-BY-4.0.txt.\par\smallskip

\noindent\textit{Editorial scope.} Counted unit: the Theorem whose source label is \texttt{thm1} in the article (source0.tex lines 677--695, source label \texttt{thm1}), delivered together with the article's own complete original proof of it (source0.tex lines 696--731, one \texttt{proof} environment of 36 lines). No mathematics is written, repaired or re-derived: statement and proof are the article's own text line for line. Required context retained uncounted: eq:ratkrylov (lines 420--422); thm:BKCS (lines 428--449). Every definition, display and label of those lines is retained unchanged. Omitted from this package and not redistributed: 1--419, 423--427, 450--676, 732--1569. Nothing inside a retained excerpt is omitted or altered: every retained body is delivered exactly as the article's lines are written, so no omission receipt of any kind accompanies this package. Citations: the retained excerpts carry no citation command at all, so no bibliography entry of the article is transcribed and no reference is redistributed. Reference adaptation: none was needed and none was made; every reference inside the delivered unit resolves inside it, so no unresolved reference or citation is produced. Printed slips kept literal: no printed word, symbol, hypothesis or number of the counted statement or of its proof was altered. Layout and class: the article's own class is \texttt{siamart190516} with packages including graphicx, amsfonts, verbatim, cite, subcaption, tikz, float, multicol, url, booktabs, mathabx, placeins, algorithm2e; the wrapper is the lane's pinned \texttt{amsart} editorial wrapper at 10pt on A4, which restates the theorem-like environments the retained text needs and adds the article's own macro definitions that the retained text needs (\texttt{\textbackslash Bignorm}, \texttt{\textbackslash R}, \texttt{\textbackslash T}, \texttt{\textbackslash Ttran}, \texttt{\textbackslash abs}, \texttt{\textbackslash bignorm}, \texttt{\textbackslash defi}, \texttt{\textbackslash diagm}, \texttt{\textbackslash norm}, \texttt{\textbackslash range}, \texttt{\textbackslash spa}, \texttt{\textbackslash tol}), with extra wrapper packages (bm, bbm, dsfont, xspace, cancel, mathtools, cleveref). The delivered numbering is the unit's own. Assets: no retained span contains a figure environment, a graphics-inclusion command, a PDF-page inclusion, a table or a TikZ picture; no image file is copied into this package, the package declares no asset dependency and no image file is redistributed with it. Licence: version arXiv:2602.20948v2 of this article is distributed under the Creative Commons Attribution 4.0 International licence, as recorded in the arXiv licence metadata for this exact version and shown on the abstract page https://arxiv.org/abs/2602.20948v2. A full rights notice is delivered at licenses/arxiv-2602.20948-CC-BY-4.0.txt. Characters: every retained excerpt is pure ASCII, so the delivered engine, XeLaTeX with its default Latin Modern text font, reports no missing character.
\begin{equation}
    \mathcal{Q}(T_{m},e_m,\Xi)  \defi \bigl\{ r(T_m)e_m \colon \text{$r$ is a rational function with poles $\Xi$} \bigr\}   \label{eq:ratkrylov}
\end{equation}

\begin{lemma}\label{thm:BKCS}
   Consider a symmetric $n\times n$
   matrix of the form
      \begin{equation*}
         T_{n}=\begin{bmatrix}
             T_{m}&\\& \star
         \end{bmatrix} + \beta \begin{bmatrix}
             &e_me_1^{\T}\\e_1e_m^{\T} &
         \end{bmatrix}\in \R^{n\times n},
     \end{equation*}
     with $T_{m} \in \R^{m \times m}$ and $\beta \in \R$.
     Let $r = p/q$ be a real rational function with poles $\Xi$ as defined above, and suppose that the eigenvalues
     of $T_m$ and $T_n$ are not zeros of $q$. Then, for any vector $v\in \R^{n}$, it holds that
     \begin{equation*}
         r(T_{n})  v \in
         \begin{bmatrix}
            r(T_{m})  \overline{v} + \mathcal{Q}(T_{m},e_m,\Xi) \\
              \R^{n-m}
         \end{bmatrix},
     \end{equation*}
     where $\overline{v}\in \R^m$ contains the first $m$ components of $v$, and $\mathcal{Q}(T_{m},e_m,\Xi)\subset \R^{m}$ denotes the rational Krylov subspace~\eqref{eq:ratkrylov}. 
\end{lemma}   

\begin{theorem}
    \label{thm1}
    With $T_{n}$ defined as in \cref{thm:BKCS}, let
    $s_{1},\dotsc,s_{m} \in \R^m$ denote orthonormal eigenvectors 
    belonging to the eigenvalues  
    $\theta_{1}\leq \dotsb\leq \theta_{m}$ of $T_{m}$, and let 
    $W_k \in \R^{n \times k}$ contain 
    orthonormal eigenvectors belonging to the smallest $k$ eigenvalues $\lambda_1 \le \cdots \le \lambda_k$ of $T_n$.
    
    Given $0<\tol_{\mathrm{ra}}<1$ and $k\leq \widehat{k}\leq m-1$ such that $\theta_k < \theta_{\widehat{k}+1}$, set $\tau=(\theta_{k}+\theta_{\widehat{k}+1})/2$ and let $r_{\tau}$ be a rational function with poles $\Xi_{\widehat{k}}$ satisfying
    \begin{equation}
        \label{con:ratapp}
        \abs{r_{\tau}(x)-\chi_{\tau}(x)} <  \tol_{\mathrm{ra}} \quad \text{for every}\quad x \in \left[\lambda_{1},\theta_{m}\right]\setminus \big(\theta_k,\theta_{\widehat{k}+1}\big).
    \end{equation}
    Let $V_{\ell}$ be an orthonormal basis of $\spa{s_{1},\dotsc,s_{\widehat{k}}}+\mathcal{Q}(T_{m},e_{m},\Xi_{\widehat{k}})$. Then
    \begin{equation*}
        \norm{(I-PP^{\Ttran})W_{k}} < 2 \cdot \tol_{\mathrm{ra}} \quad {\mathrm{with}} \quad P=\diagm(V_{\ell},I).
    \end{equation*}
\end{theorem}

\begin{proof} 
    It suffices to show that $\norm{(I-PP^{\Ttran})w}\leq 2 \cdot \tol_{\mathrm{ra}}$ for every $w = W_{k} \alpha$ with $\alpha = [\alpha_1,\ldots,\alpha_k]^\Ttran$ such that $\norm{\alpha}=1$.
    By Cauchy interlacing, $\lambda_{1}\leq \lambda_{i}\leq \theta_{k} < \tau$ for $i = 1,\ldots, k$.
    Using the error bound~\cref{con:ratapp}, we have that 
    \begin{equation*}
        \norm{r_{\tau}(T_{n})w-w}^{2} =
        \Bignorm{\sum_{i=1}^{k}\alpha_{i}\bigl(r_{\tau}(\lambda_i)-1\bigr) w_{i}}^{2} = \sum_{i=1}^{k} \alpha_{i}^{2}\bigl(r_{\tau}(\lambda_{i})-\chi_{\tau}(\lambda_{i})\bigr)^{2} <  \tol_{\mathrm{ra}}^{2}. 
    \end{equation*} 
    This implies $\norm{(I-PP^{\Ttran})w}< \norm{(I-PP^{\Ttran})r_{\tau}(T_{n})w}+\tol_{\mathrm{ra}}$. It remains to show that the first term of the sum is also bounded 
    by $\tol_{\mathrm{ra}}$. For this purpose, we let 
    $\overline{w}\in\R^{m}$ contain the first $m$ components of $w$ and 
    recall from~\cref{thm:BKCS} that
    \begin{equation*}
        r_{\tau}(T_{n})w - \begin{bmatrix}
            r_{\tau}(T_{m})\overline{w}\\0
        \end{bmatrix}\in 
            \range\begin{bmatrix}
            V_{\ell}& 0 \\ 0 & I
        \end{bmatrix}=\range(P).
    \end{equation*}
    In turn,
    \begin{equation*}
        \begin{aligned}
            \norm{(I-PP^{\Ttran})r_{\tau}(T_{n})w} &= \bignorm{(I-V_{\ell}V_{\ell}^{\Ttran})r_{\tau}(T_{m})\overline{w}} =\Bignorm{(I-V_{\ell}V_{\ell}^{\Ttran})\sum_{i=1}^{m}r_{\tau}(\theta_{i})(s_{i}^{\Ttran}\overline{w})s_{i}} \\ 
            &= \Bignorm{(I-V_{\ell}V_{\ell}^{\Ttran})\sum_{i=\widehat{k}+1}^{m}r_{\tau}(\theta_{i})(s_{i}^{\Ttran}\overline{w})s_{i}}
            \leq \Bignorm{\sum_{i=\widehat{k}+1}^{m}r_{\tau}(\theta_{i})(s_{i}^{\Ttran}\overline{w})s_{i}},
        \end{aligned}
    \end{equation*}
    where we inserted the spectral decomposition $T_{m}=\sum_{i=1}^{m}\theta_{i}s_{i}s_{i}^{\Ttran}$ in the second equality and used $s_{i}\in\range(V_{\ell})$ for $1\leq i\leq \widehat{k}$ in the third equality. The proof is completed by  
    \begin{equation*}
        \Bignorm{\sum_{i=\widehat{k}+1}^{m}r_{\tau}(\theta_{i})(s_{i}^{\Ttran}\overline{w})s_{i}}^{2}
        =\sum_{i=\widehat{k}+1}^{m}r_{\tau}^{2}(\theta_{i})(s_{i}^{\Ttran}\overline{w})^{2} \leq \norm{\overline{w}}^{2}\max_{\widehat{k}+1\leq i\leq m}\abs{r_{\tau}(\theta_{i})}^{2}
        < \tol_{\mathrm{ra}}^{2},  
    \end{equation*}
    where we applied the error bound~\cref{con:ratapp} once more. 
\end{proof}

\end{document}
